\subsection{Euler characteristics and the genus-zero ruling}
\label{subsec:equality-torsor-euler}

The group action in Corollary~\ref{cor:normalized-equality-torsor}
gives a global filtration of the normalized cover's direct-image algebra.
Its quotients determine the Euler characteristic. In the genus-zero
ruling case, this also gives a lower bound for the defect of rational
singularities and an obstruction to extending the distinguished ruling.
The canonical-section argument in the following subsection turns these
defects into exclusions and a characteristic-three residual case.

Degree filtrations from affine coordinate changes are standard for
additive torsors; see the $\alpha_{L,a}$-torsor construction in
\cite[equation~(0.2.17), pp.~38--40]{CDL}. The group used here has a
deformed multiplicative coproduct. We give the descent calculation
explicitly, including its torsion character twist.

\begin{proposition}\label{prop:deformed-torsor-degree-filtration}
Let $X$ be a scheme of characteristic $p>0$, let $L$ be an invertible
$\OO_X$-module, and let $s\in H^0(X,L)$. Define $G_s$ as in
Corollary~\ref{cor:smooth-torsor}: if $s_j=g_{ij}^{-1}s_i$ in two
frames of $L$, its local Hopf algebras are
\[
 \OO_X[u_i]/(u_i^p),\qquad
 \Delta(u_i)=u_i\otimes1+1\otimes u_i+s_i u_i\otimes u_i,
 \qquad u_j=g_{ij}u_i.
\]
For every finite flat $G_s$-torsor $\pi:Q\to X$, there is a filtration
by finite locally free submodules
\[
 0=F_{-1}\subset F_0=\OO_X\subset F_1\subset\cdots
 \subset F_{p-1}=\pi_*\OO_Q
\]
with $\operatorname{rank}F_i=i+1$. There is an invertible module $P$ such
that
\begin{equation}\label{eq:deformed-torsor-graded}
 F_i/F_{i-1}\simeq L^{\otimes i}\otimes P^{\otimes i}
 \quad(0\leq i<p),\qquad
 P^{\otimes p}\simeq\OO_X,\qquad
 P|_{Z(s)}\simeq\OO_{Z(s)}.
\end{equation}
Here $Z(s)$ is the scheme-theoretic zero subscheme.
The filtration is intrinsic to the given torsor action.
\end{proposition}

\begin{proof}
The gluing rule identifies the underlying module of the group algebra
with
\[
 \mathcal A=(G_s\to X)_*\OO_{G_s}
 =\bigoplus_{j=0}^{p-1}L^{\otimes j}.
\]
Thus its degree submodules
$V_i=\bigoplus_{j=0}^iL^{\otimes j}$ are global subbundles, with
$V_i/V_{i-1}\simeq L^{\otimes i}$.
The local operation $x*y=x+y+s_i xy$ has inverse
$-x/(1+s_i x)$. The denominator is a unit because $x^p=0$.
The formula
\[
 \vartheta:G_s\longrightarrow\mu_p,\qquad
 \vartheta(u_i)=1+s_i u_i
\]
defines a global character: it is unchanged on overlaps,
is multiplicative for the group operation, and has $p$th power one.

In a fixed frame of $L$, translation by a section $v$ sends
\[
 u\longmapsto v+(1+sv)u.
\]
This substitution preserves every degree submodule $V_i$.
On its degree-$i$ quotient the substitution is multiplication by
$(1+sv)^i=\vartheta(v)^i$, a unit. The inverse translation gives
the inverse linear map. These assertions hold over arbitrary
base rings, including nonreduced ones.

Take an fpqc cover $U\to X$ on which the torsor has a section;
one may take $U=Q$. The section identifies $Q_U$ with $(G_s)_U$,
and the descent datum on
$(\pi_*\OO_Q)|_U\simeq\mathcal A|_U$ is given by translations.
It therefore preserves all the $V_i|_U$, and the cocycle condition
restricts to these submodules. Effective descent and full faithfulness
for quasi-coherent modules give their descents $F_i$, the inclusions
between them, and their maps to $\pi_*\OO_Q$
\cite[Proposition~35.5.2, Tag~023T]{Stacks}.
Injectivity and exactness can be checked after faithfully flat pullback.
Consequently the quotients pull back to $V_i/V_{i-1}$.
The modules $F_i$ and all these quotients are finite locally free,
because that property descends under an fpqc cover
\cite[Lemma~35.7.6, Tag~05B2]{Stacks}.
In particular the construction has no exceptional finite-length
quotient at a point of $X$.

After factoring out the ordinary $L$-transition, the degree-one
quotient has the scalar descent cocycle $\vartheta(v)$.
Let $P$ be its descended line bundle; equivalently,
$P=(F_1/F_0)\otimes L^{-1}$.
The leading coefficient in degree $i$ is the $i$th power of the
degree-one coefficient. This gives the first isomorphism in
\eqref{eq:deformed-torsor-graded}. The scalar cocycle has $p$th
power one, so $P^{\otimes p}\simeq\OO_X$.
On $Z(s)$ the scalar cocycle is itself one, giving the last
isomorphism in \eqref{eq:deformed-torsor-graded}, including when
$Z(s)$ is nonreduced. Constants are fixed by every translation,
so $F_0$ is the unit submodule $\OO_X$.
Changing the chosen torsor section is another translation,
which preserves the flag. The resulting filtration is therefore
independent of these choices.
\end{proof}

There is also an exact identification of the torsion twist.
Suppose $Q\to T$ is equivariant for $\vartheta:G_s\to\mu_p$,
where $T\to X$ is a $\mu_p$-torsor. Let $W$ be the weight-one
line in $(T\to X)_*\OO_T$, defined locally by the coordinate
$z$ with coaction $z\mapsto z\otimes t$.
A section of $Q$ gives a section of $T$. If two sections of $Q$
differ by $v$, the induced sections of $T$ differ by $\vartheta(v)$.
The weight-one coordinate therefore has the same scalar cocycle
as $P$, and $P\simeq W$ by descent.
For the equality torsor $N\to X$, the normalization map $N\to T$
has this equivariance by the coaction constructed in
Corollary~\ref{cor:normalized-equality-torsor}.
Thus its graded pieces are exactly
$\OO_X(iaY)\otimes W^{\otimes i}$.
Defining $W$ by its weight fixes the convention independently of
whether the original Kummer algebra is written using $M$ or $M^{-1}$.

\begin{corollary}\label{cor:normalized-equality-euler}
Let $X$ be a smooth projective integral surface over an algebraically
closed field of characteristic $p>0$. Let $Y\subset X$ be a smooth
ample integral curve, and let $M\ne\OO_X$ have order $p$ with
$M|_Y\simeq\OO_Y$. Let $\pi:N\to X$ be the normalization of the
$\mu_p$-torsor defined by a chosen trivialization of $M^p$.
Let $a\geq1$ be the conductor integer of
Proposition~\ref{prop:normalized-mup-conductor}, characterized by
the order $pa$ of the logarithmic differential along $Y$.
If $2g(Y)-2=apY^2$, then
\begin{equation}\label{eq:normalized-equality-euler}
 \chi(\OO_N)=p\chi(\OO_X)
 -\frac{p(p-1)a\bigl((p+1)a-3\bigr)}{12}Y^2.
\end{equation}
In particular, if $a=1$, then
\begin{equation}\label{eq:normalized-smallest-equality-euler}
 \chi(\OO_N)=p\chi(\OO_X)
 -\frac{p(p-1)(p-2)}{12}Y^2.
\end{equation}
\end{corollary}

\begin{proof}
Corollary~\ref{cor:normalized-equality-torsor} makes $\pi:N\to X$
a finite flat torsor for $G_{s_a}$, with
$L=\OO_X(aY)$ and $s_a$ its canonical section.
Proposition~\ref{prop:deformed-torsor-degree-filtration}
therefore applies. Its line bundle $P$ has $P^p\simeq\OO_X$ and
$P|_Y\simeq\OO_Y$, since $Y\subset Z(s_a)=aY$.

For the line bundle $P$ on the proper scheme $X$, the function
$n\mapsto\chi(P^{\otimes n})$ is a numerical polynomial over
$\mathbf Q$ \cite[Lemma~33.45.1, Tag~0BEM]{Stacks}.
Because $P^p\simeq\OO_X$, it is periodic with period $p$, so it is
constant: the difference of a
nonconstant polynomial at $n+p$ and $n$ has a nonzero leading
coefficient. Hence
\[
 \chi(P^{\otimes i})=\chi(\OO_X)
 \qquad\text{for every integer }i.
\]
Put $d=Y^2$. For $j\geq1$, the Cartier divisor sequence gives
\[
 0\longrightarrow P^{\otimes i}((j-1)Y)
 \longrightarrow P^{\otimes i}(jY)
 \longrightarrow\OO_Y(jY)\longrightarrow0.
\]
The degree formula on the smooth proper integral curve $Y$ gives
\[
 \chi(\OO_Y(jY))=jd+1-g(Y)
\]
\cite[Definition~33.44.1 and Lemmas~33.44.5 and~33.44.7,
Tag~0AYQ]{Stacks}.
Summing these sequences yields
\[
 \chi\bigl(P^{\otimes i}(nY)\bigr)
 =\chi(\OO_X)+\frac{n(n+1)}2d+n(1-g(Y))
 \qquad(n\geq0).
\]
Set $n=ia$ and use $1-g(Y)=-apd/2$. The $i$th graded quotient
of the filtration has Euler characteristic
\[
 \chi(\OO_X)+\frac d2\bigl(a^2i^2-a(ap-1)i\bigr).
\]
Since $\pi$ is finite, cohomology of $\OO_N$ equals cohomology of
$\pi_*\OO_N$: the affine \v{C}ech complexes on $X$ and their inverse
images on $N$ agree. Additivity along the filtration and
\[
 \sum_{i=0}^{p-1}i=\frac{p(p-1)}2,\qquad
 \sum_{i=0}^{p-1}i^2=\frac{p(p-1)(2p-1)}6
\]
now give
\[
 \begin{split}
 \chi(\OO_N)
 &=p\chi(\OO_X)
 +\frac d2\left(
 a^2\frac{p(p-1)(2p-1)}6
 -a(ap-1)\frac{p(p-1)}2\right)\\
 &=p\chi(\OO_X)
 -\frac{p(p-1)a((p+1)a-3)}{12}d.
 \end{split}
\]
Taking $a=1$ proves the second formula.
\end{proof}

\begin{proposition}\label{prop:genus-zero-equality-defect}
Let $X$ be a smooth projective integral surface over an algebraically
closed field $k$ of characteristic $p>0$.
Let $Y\subset X$ be a smooth integral very ample divisor, and let
$M\ne\OO_X$ have order $p$ with $M|_Y\simeq\OO_Y$.
Suppose $2g(Y)-2=pY^2$, and let $\pi:N\to X$ be the normalization
of its chosen $\mu_p$-torsor.
Choose a ruling presentation
$k(N)=k(C)(t)$ as in
Corollary~\ref{cor:very-ample-equality-almost-ruled}, and suppose
$g(C)=0$.
Let $r:S\to N$ be the proper birational morphism from a point-blowup
model of $C\times\PP^1$ supplied by
Proposition~\ref{prop:distinguished-ruled-cover}. Put
\[
 \mathcal T=R^1r_*\OO_S,\qquad
 \epsilon=\dim_kH^0(N,\mathcal T),\qquad d=Y^2.
\]
Then $\mathcal T$ has finite support and
\begin{equation}\label{eq:genus-zero-equality-defect}
 H^1(N,\OO_N)=0,\qquad
 H^2(N,\OO_N)\simeq H^0(N,\mathcal T),\qquad
 \chi(\OO_N)=1+\epsilon.
\end{equation}
If $p\ne3$, then $\epsilon\equiv-1\pmod p$ and
$\epsilon\geq p-1$. In particular $N$ has a nonrational singularity
outside $D=\pi^*Y$, and the distinguished rational map
$N\dashrightarrow C$ does not extend to a morphism.

If $p=3$, write $d=2e$. Then $\epsilon\equiv-e-1\pmod3$.
Rational singularities, or extension of that rational map, can occur
only if $d\equiv4\pmod6$.
\end{proposition}

\begin{proof}
The field is algebraically closed, so $C\simeq\PP^1$.
A point blowup of a smooth surface has pushforward $\OO$ and
vanishing positive higher direct images
\cite[Lemma~54.3.4, Tag~0AGP]{Stacks}.
The vanishing also follows directly from its two affine charts:
on their intersection, every Laurent monomial in the chart parameter
belongs to one chart ring or the other, so the degree-one
\v{C}ech differential is surjective.
The two-chart computation for $\PP^1_R$ has the same conclusion.
Applying Leray to the point-blowup presentation of $S$ gives
\[
 H^0(S,\OO_S)=k,\qquad
 H^1(S,\OO_S)=H^2(S,\OO_S)=0.
\]

Normality gives $r_*\OO_S=\OO_N$. Indeed proper direct-image
coherence \cite[III, Theorem~3.2.1]{EGA} makes the former a
finite algebra, and its sections over an affine open lie in the
common function field and are integral over the normal base ring.
The proper birational map is an isomorphism near every
codimension-one point of $N$ \cite[Tag~0BFP]{Stacks}.
Thus the coherent sheaf $\mathcal T$ has finite support.

We also have $R^ir_*\OO_S=0$ for $i\geq2$.
To verify this without a smooth-fibre assertion, use a projective
embedding of $S$ over $N$; it exists by the closed graph construction.
Every fibre of $r$ has dimension at most one.
At a closed point of $N$, choose a relative hyperplane containing
no component of its fibre, then another avoiding the finite
intersection of the first with that fibre.
These choices are possible because $k$ is infinite.
Properness allows an affine neighbourhood on which the two
hyperplanes have empty common intersection with $S$.
Their complements give an affine cover with affine intersection,
so its \v{C}ech cohomology vanishes in degrees at least two.
These neighbourhoods prove the higher-direct-image assertion.

The low-degree Leray sequence, using
$H^1(S,\OO_S)=H^2(S,\OO_S)=0$, gives the two cohomology assertions
in \eqref{eq:genus-zero-equality-defect}. Since $\mathcal T$
has finite support, its Euler characteristic is $\epsilon$,
and Leray also gives
$1=\chi(\OO_S)=\chi(\OO_N)-\epsilon$.
A nonzero stalk of $\mathcal T$ witnesses a nonrational singularity:
the local ring is a two-dimensional normal Nagata domain and its
localized smooth model is a normal modification with nonzero
$H^1$ \cite[Definition~54.8.3, Tag~0B4N]{Stacks}.
The equality conclusion of
Proposition~\ref{prop:normalized-mup-conductor}
makes $N$ regular near $D$, so this singularity lies outside $D$.

For $p=2$, \eqref{eq:normalized-smallest-equality-euler} gives
$\chi(\OO_N)=2\chi(\OO_X)$.
For a prime $p\geq5$, the equality $2g(Y)-2=pd$ makes $d=2e$
even, and $(p-1)(p-2)$ is divisible by six.
The correction term in
\eqref{eq:normalized-smallest-equality-euler} is therefore
$p(p-1)(p-2)e/6$, an integer divisible by $p$.
For every $p\ne3$, it follows that
$1+\epsilon=\chi(\OO_N)\equiv0\pmod p$.
Nonnegativity of $\epsilon$ yields $\epsilon\geq p-1$.
For $p=3$, the same Euler formula is
$\chi(\OO_N)=3\chi(\OO_X)-e$, giving
$\epsilon\equiv-e-1\pmod3$.
Thus $\epsilon=0$ requires $e\equiv2\pmod3$,
or $d\equiv4\pmod6$.

Finally suppose the distinguished rational map extends to
$h:N\to C$. Let $q:S\to C$ be the ruling morphism.
Then $hr=q$, since they agree generically and $C$ is separated.
The blowup and projective-line calculations give
$q_*\OO_S=\OO_C$ and $R^iq_*\OO_S=0$ for $i>0$.
The two-hyperplane argument, now for the projective map $h$
with one-dimensional fibres, gives $R^ih_*\OO_N=0$ for $i\geq2$.
The low-degree Leray sequence for $q=hr$ is
\[
 0\longrightarrow R^1h_*\OO_N
 \longrightarrow R^1q_*\OO_S
 \longrightarrow h_*\mathcal T
 \longrightarrow R^2h_*\OO_N.
\]
Hence $R^1h_*\OO_N=0$ and $h_*\mathcal T=0$.
Pushforward of a finite-support coherent sheaf is faithful here:
its scheme-theoretic support is finite over $k$, hence over $C$,
and pushforward is restriction of scalars on affine opens.
Consequently $\mathcal T=0$.
Also $h_*\OO_N=q_*\OO_S=\OO_C$, so $\chi(\OO_N)=1$.
This contradicts the congruence when $p\ne3$.
In characteristic three it requires the stated congruence on $d$.
\end{proof}

For $p=3$, the same calculation gives $\epsilon\geq2$ when
$d\equiv0\pmod6$ and $\epsilon\geq1$ when $d\equiv2\pmod6$.
The remaining residue class is a necessary numerical possibility;
no example is asserted.

Whenever a distinguished ruling $N\dashrightarrow C$ extends,
the last cohomological argument works for any genus and gives
$\chi(\OO_N)=1-g(C)$.
In the smallest equality case,
\eqref{eq:normalized-smallest-equality-euler} consequently gives
\[
 g(C)\equiv1\pmod p\quad(p\ne3),\qquad
 \frac{Y^2}{2}\equiv g(C)-1\pmod3\quad(p=3).
\]
The extension over a positive-genus base follows from
Proposition~\ref{prop:positive-genus-picard-fibration}.
For a genus-zero base the preceding proposition instead requires
indeterminacy when $p\ne3$.
For a very ample divisor, the canonical subsystem in
Corollary~\ref{lem:divided-differential-subsystem} supplies a further
positivity constraint, and Theorem~\ref{thm:strict-picard-all}
excludes the smallest equality in every characteristic.
The merely ample and higher-conductor questions remain at the
scope described in the final section.
